\documentclass[a4paper]{article}

%%packages%%
\usepackage[french]{babel}
\usepackage{amsmath,amssymb}
\usepackage{parskip, setspace}
\usepackage{geometry}
\usepackage{bbold}
\usepackage{mathtools, stmaryrd}
\usepackage{graphicx}
\usepackage{float}
\usepackage{tikz}
\usepackage{pgfplots}
\pgfplotsset{compat=1.18}
\usepackage{hyperref}


%%format page%%
\author{Mohammed JABRI , Bastien MAGGI \\
Encadrant : Benjamin CUER}
\date{}
\title{Projet M1 SSD:\\
Détection sérologique chez la faune sauvage : estimation de seuils de
positivité par modèles de mélange}
\geometry{margin=2.5cm}
\onehalfspacing

\begin{document}
\maketitle
\tableofcontents
\newpage

\section{Introduction}

\subsection{Contexte}

La surveillance des maladies émergentes est essentielle pour détecter précocement la circulation d’agents pathogènes et mieux comprendre les risques de transmission entre les animaux et l’être humain. De nombreuses maladies infectieuses émergentes ont une origine zoonotique, c’est-à-dire qu’elles proviennent d’un réservoir animal avant de franchir la barrière d’espèce. La faune sauvage joue donc un rôle important dans cette surveillance, car certaines espèces peuvent héberger des virus ou d’autres agents infectieux sans présenter de signes cliniques visibles. L’étude sérologique de ces populations permet d’identifier les individus ayant déjà été exposés à un pathogène grâce à la détection d’anticorps qui sont des protéines produites par le système immunitaire en réponse à un antigène, comme une protéine virale. Leur détection dans le sérum permet donc d’indiquer qu’un individu a été exposé au pathogène étudié, même lorsque l’agent infectieux n’est plus directement présent. Cette approche est particulièrement utile chez les animaux sauvages, pour lesquels la détection directe par PCR. La PCR, ou réaction de polymérisation en chaîne, est une méthode permettant d’amplifier et de détecter le matériel génétique d’un agent pathogène. Elle détecte surtout les infections actives peut être difficile en raison de la courte durée de présence du virus dans l’organisme. Ainsi, la sérologie constitue un outil important pour estimer l’exposition passée, suivre la circulation des pathogènes et mieux évaluer les risques d’émergence à partir de la faune sauvage. \\

Plusieurs méthodes existent pour doser les anticorps d'un virus. La méthode classique est le test ELISA (Enzyme-Linked Immunosorbent Assay). Dans cette technique, des antigènes viraux sont fixés au fond d’une microplaque. Le sérum sanguin de l’animal étudié est ajouté ; si des anticorps spécifiques sont présents, ils se lient aux antigènes. Un anticorps secondaire couplé à une enzyme est ensuite ajouté, produisant une réaction colorée proportionnelle à la quantité d’anticorps présents. L’intensité de cette réaction est mesurée par un lecteur de plaque et fournit une valeur quantitative continue des anticorps. \\

Une alternative plus récente est la technologie Luminex multiplex basée sur des microsphères fluorescentes. Chaque bille est associée à un antigène spécifique, permettant la détection simultanée d’anticorps contre plusieurs pathogènes. Après liaison des anticorps présents dans le sérum, un anticorps secondaire fluorescent est ajouté et la fluorescence est mesurée par un système laser. Le résultat obtenu est une valeur appelée Mean Fluorescence Intensity (MFI), proportionnelle à la concentration relative d’anticorps. \\

Ces valeurs MFI constituent la base des analyses statistiques ultérieures visant à distinguer individus séroné-
gatifs et séropositifs.

Cependant, l’interprétation des valeurs MFI n’est pas toujours simple, en particulier dans les études portant sur la faune sauvage. Contrairement à de nombreuses applications humaines ou vétérinaires bien établies, les espèces sauvages sont souvent des espèces non modèles, pour lesquelles les tests sérologiques sont peu validés et les échantillons de contrôle positifs ou négatifs sont rarement disponibles. Il devient alors difficile de définir un seuil fiable permettant de distinguer les individus séronégatifs des individus séropositifs. De plus, les valeurs MFI sont des mesures continues : elles ne donnent pas directement un statut sérologique, mais doivent être interprétées à l’aide d’une méthode statistique. Cette difficulté est importante dans le cadre de la surveillance des maladies émergentes, car une mauvaise classification des individus peut entraîner une mauvaise estimation de la séroprévalence et donc une mauvaise compréhension de la circulation d’un pathogène dans une population sauvage. Ces limites ont été soulignées dans plusieurs travaux portant sur l’analyse sérologique en écologie et en surveillance virale, notamment pour les méthodes ELISA et Luminex.

Plusieurs approches statistiques peuvent être utilisées pour analyser ces données \cite{kain2025} \cite{peeters2023}. \\
Une méthode simple consiste à définir un seuil de positivité, appelé cut-off, à partir des valeurs supposées négatives. Cette approche est facile à appliquer, mais elle dépend fortement du choix du seuil. D’autres méthodes, comme les modèles de mélange gaussien, cherchent plutôt à représenter les données comme provenant de plusieurs groupes sous-jacents, par exemple un groupe d’individus supposés séronégatifs et un groupe d’individus supposés séropositifs. Dans ce projet, nous nous limiterons donc à l’étude de deux approches : la Méthode des valeurs extrêmes (outliers) et le modèle de mélange gaussien à deux composantes.

La problématique de ce travail est donc la suivante : comment analyser des valeurs MFI issues d’une étude sérologique en faune sauvage afin de distinguer les individus supposés séronégatifs des individus supposés séropositifs ?



\textbf{Mean Fluorescence Intensity (MFI)} : est une variable continue mesurant l'intensité moyenne de fluorescence d'un anticorps. C'est une mesure indirecte de la quantité d'un anticorps spécifique. C'est une variable continue et positive mesurée en laboratoire par un appareil.\\

\textbf{ELISA ou LUMINEX} : Les tests sérologiques tels que ELISA ou Luminex produisent des mesures quantitatives continues d’anticorps (MFI). Afin d’interpréter biologiquement ces valeurs, une étape statistique est nécessaire pour classer les individus en séronégatifs ou séropositifs, généralement via une méthode de cut-off. Cette classification permet ensuite d’estimer la séroprévalence au sein de la population étudiée.

\textbf{Cut-off} : correspond à un seuil statistique permettant de classer les individus comme séronégatifs ou séropositifs à partir de mesures continues d’anticorps telles que les valeurs MFI.

\subsection{Objectifs}

L’objectif de ce projet est d’étudier l’utilisation d’un modèle de mélange gaussien à deux composantes pour l’analyse de données sérologiques. Ce modèle permet de représenter les mesures observées comme provenant de deux populations sous-jacentes : une population supposée séronégative et une population supposée séropositive. \\
Le but est ensuite de comparer cette approche à une méthode plus classique de détermination d’un seuil de positivité, appelée méthode du cut-off. Ces deux méthodes seront appliquées à des données sérologiques issues de chauves-souris afin d’estimer la proportion d’individus ayant potentiellement déjà été en contact avec un virus émergent de type Ebola.

Dans ce travail, nous nous concentrerons donc sur deux approches :
\begin{itemize}
    \item Méthode des valeurs extrêmes (outliers)
    \item Modèle de mélange Gaussien (GMM) avec 2 classes
\end{itemize}

\section{Présentation des méthodes}

Dans cette section, nous présentons les deux méthodes statistiques étudiées dans ce projet pour analyser les valeurs MFI issues des données sérologiques. La première approche est la méthode des valeurs extrêmes (outliers), qui consiste à définir un seuil de positivité. La seconde approche est le modèle de mélange gaussien à deux composantes, qui permet de représenter les données comme provenant de deux groupes sous-jacents. Ces deux méthodes seront ensuite appliquées et comparées dans la partie pratique.

\subsection{Méthode des valeurs extrêmes (outliers)}

\subsubsection{Principe général} 

La méthode du cut-off utilisée dans ce projet consiste à définir un seuil de positivité à partir des valeurs les plus faibles des données MFI, supposées représenter les individus séronégatifs. Une fois ce seuil calculé, les observations ayant une valeur MFI inférieure ou égale au seuil sont classées comme supposées séronégatives, tandis que les observations ayant une valeur MFI supérieure au seuil sont classées comme supposées séropositives.

Dans notre application, cette méthode repose donc sur l'hypothèse que les faibles valeurs MFI correspondent aux individus ne présentant pas de réponse sérologique détectable contre le pathogène étudié. Les valeurs situées au-delà du seuil défini sont alors interprétées comme des réponses sérologiques plus élevées, pouvant indiquer la présence d'anticorps.

\subsubsection{Modélisation statistique} 

Soit \(X = (X_1, \ldots, X_n)\) un échantillon de variables aléatoires indépendantes représentant les mesures de MFI. L’indice \(i\) désigne le numéro de l’observation, avec \(i = 1, \ldots, n\). Autrement dit, \(X_i\) correspond à la valeur MFI mesurée pour le \(i\)-ème individu.

On suppose que les observations \(X_i\) issues de la population séronégative suivent approximativement une loi normale :

\[
X_i \sim \mathcal{N}(\mu, \sigma^2)
\]

où \(\mu\) et \(\sigma^2\) désignent respectivement la moyenne et la variance de la population séronégative.

Dans ce cadre, les observations issues de la population séropositive sont considérées comme des valeurs extrêmes situées au-delà d'un certain seuil à définir, aussi appelé \textit{cut-off}. \\

\textbf{Définition du seuil} \\

Le seuil sérologique est défini comme :

\[
c = \mu + k \sigma
\]

où $k$ est un paramètre fixé. 



Dans la pratique, $k = 3$ est fréquemment utilisé, ce qui correspond à un seuil tel que :

\[
c = \mu + 3\sigma
\]

%Sous l’hypothèse de normalité, environ 99,7\% des observations issues de la population négative sont inférieures à ce seuil. \\
Dans le cas d’une loi normale, le seuil \(c=\mu+3\sigma\) correspond à une valeur située trois écarts-types au-dessus de la moyenne. La probabilité qu’une observation issue de la population négative dépasse ce seuil est très faible : elle est d’environ \(0{,}135\%\). Les valeurs situées au-delà de ce cut-off peuvent donc être considérées comme des valeurs extrêmes. 

\newpage

La figure suivante illustre le positionnement du seuil \(c=\mu+3\sigma\) sur une loi normale. La partie située à droite de ce seuil correspond à une proportion très faible des observations, ce qui justifie l’interprétation de ces valeurs comme des valeurs extrêmes.

\begin{figure}[h!]
\centering
\begin{tikzpicture}
\begin{axis}[
    no markers,
    domain=-4.5:4.5,
    samples=200,
    axis lines=left,
    xlabel={Valeur standardisée},
    ylabel={Densité},
    width=11cm,
    height=6cm,
    ymin=0,
    ymax=0.45,
    xtick={0,3},
    xticklabels={$\mu$, $\mu+3\sigma$},
    ytick=\empty,
    enlargelimits=false
]

% Courbe normale
\addplot[thick, blue] {1/sqrt(2*pi)*exp(-x^2/2)};

% Ligne verticale à mu
\addplot[dashed] coordinates {(0,0) (0,0.399)};

% Ligne verticale du cutoff
\addplot[red, thick] coordinates {(3,0) (3,0.02)};
\node[red, above] at (axis cs:3,0.03) {cut-off};

% Zone de queue à droite de 3
\addplot[
    domain=3:4.5,
    samples=100,
    fill=red!20,
    draw=none
] {1/sqrt(2*pi)*exp(-x^2/2)} \closedcycle;

\end{axis}
\end{tikzpicture}
\caption{Illustration d’une loi normale et du seuil \(c=\mu+3\sigma\). La zone située à droite du cut-off représente les observations extrêmes, soit environ \(0{,}135\%\) des données dans le cas d’une loi normale.}
\end{figure}

\textbf{Estimation des paramètres} \\

En l’absence de données de contrôle négatif, [Les contrôles négatifs sont des échantillons dont on sait qu’ils proviennent d’individus non exposés au pathogène. Ils servent à estimer ce que donne le test chez des individus réellement séronégatifs], les paramètres $\mu$ et $\sigma$ sont estimés à partir des données observées :

\[
\hat{\mu} = \frac{1}{n} \sum_{i=1}^{n} X_i, \quad
\hat{\sigma} = \sqrt{\frac{1}{n-1} \sum_{i=1}^{n} (X_i - \hat{\mu})^2}
\]

Le seuil estimé est alors donné par :

\[
\hat{c} = \hat{\mu} + 3 \hat{\sigma}
\]

\textbf{Classification des individus} \\

Chaque individu est ensuite classé selon la règle :

\[
\text{Classe}_i =
\begin{cases}
1 & \text{si } X_i > \hat{c} \\
0 & \text{sinon}
\end{cases}
\]

où $1$ correspond à un individu séropositif et $0$ à un individu séronégatif. \\

\newpage 

\subsubsection{Estimation de la prévalence}

La prévalence sérologique est estimée comme la proportion d'individus classés positifs :

\[
p = \frac{1}{n} \sum_{i=1}^{n} \mathbb{1}_{\{X_i > \hat{c}\}}
\]

où $\mathbb{1}$ désigne la fonction indicatrice. \\


\textbf{Limites de la méthode} \\

Bien que simple à mettre en œuvre, cette approche repose sur plusieurs hypothèses fortes :

\begin{itemize}
    \item la distribution des valeurs négatives est supposée gaussienne ;
    \item la population négative est supposée majoritaire ;
    \item la présence de valeurs positives peut biaiser l’estimation de $\mu$ et $\sigma$.
\end{itemize}

Dans nos données, certaines distributions présentent une forte asymétrie ainsi qu’un chevauchement important entre populations négatives et positives. Dans ces situations, l’hypothèse de normalité peut être discutable et conduire à des seuils moins robustes. De plus, le choix de la proportion d’individus considérés comme contrôles influence directement l’estimation du cut-off.



\subsection{Modèle de mélange Gaussien (GMM)}

\subsubsection{Principe général}

Les modèles de mélange sont des modèles statistiques utilisés lorsque les données observées peuvent provenir de plusieurs populations sous-jacentes non directement connues. Ils permettent de représenter une distribution globale comme une combinaison de plusieurs distributions plus simples, appelées composantes du mélange. Chaque composante représente alors un groupe latent possible dans les données.
Les modèles de mélanges gaussiens sont fréquemment utilisés
pour l’analyse des données sérologiques \cite{garnier2017}.

Lorsque chaque composante du mélange est supposée suivre une loi normale, on parle de modèle de mélange gaussien, ou \textit{Gaussian Mixture Model} GMM. Dans notre cas, on considère un mélange de deux composantes : une composante associée aux individus supposés séronégatifs, qui présentent généralement des valeurs MFI faibles, et une composante associée aux individus supposés séropositifs, qui présentent généralement des valeurs MFI plus élevées.

Les GMM sont utiles car ils permettent à la fois de modéliser des distributions plus flexibles qu’une seule loi normale et de classer les observations selon leur probabilité d’appartenance à chaque composante. Dans le cadre de l’analyse sérologique, cette approche permet donc d’attribuer à chaque individu une probabilité d’appartenir au groupe séronégatif ou au groupe séropositif. Une référence classique pour approfondir ces modèles est l’ouvrage de McLachlan et Peel, \textit{Finite Mixture Models}.

\newpage 

\subsubsection{Modélisation statistique}

Soit \(X = (X_1, \ldots, X_n)\) un échantillon de variables aléatoires indépendantes et identiquement distribuées, notées i.i.d. Chaque \(X_i\) représente la valeur MFI mesurée pour le \(i\)-ème individu, avec \(i = 1,\ldots,n\).\\
Le modèle de mélange gaussien à deux composantes s'écrit d'abord sous la forme :

\[
f(x) = \pi_1 \, \phi(x;\mu_1,\sigma_1) + \pi_2 \, \phi(x;\mu_2,\sigma_2)
\]

où \(\pi_1\) et \(\pi_2\) sont les proportions associées aux deux composantes du mélange. Ces proportions vérifient :

\[
\pi_1 + \pi_2 = 1
\]

avec \(\pi_1 \in (0,1)\) et \(\pi_2 \in (0,1)\).

Ainsi, on peut poser \(\pi_1 = \pi\) et \(\pi_2 = 1-\pi\). Le modèle devient alors :

\[
f(x) = \pi \, \phi(x;\mu_1,\sigma_1) + (1-\pi) \, \phi(x;\mu_2,\sigma_2)
\]

où :
\begin{itemize}
    \item \(\phi(x;\mu,\sigma)\) désigne la densité d'une loi normale \(\mathcal{N}(\mu,\sigma^2)\) :

    \[
    \phi(x;\mu,\sigma) =
    \frac{1}{\sqrt{2\pi}\sigma}
    \exp\left(
    -\frac{(x-\mu)^2}{2\sigma^2}
    \right)
    \]
    \item $(\mu_1,\sigma_1)$ et $(\mu_2,\sigma_2)$ sont les paramètres de la gaussienne de chaque composante. \\
\end{itemize}

\textbf{Variables latentes} \\

On introduit une variable latente $Z_i \in \{1,2\}$ indiquant l’appartenance de l’observation $X_i$ à l’une des deux composantes du mélange. Dans notre cas, cette variable permet d’interpréter chaque observation comme appartenant soit à la composante associée aux individus supposés séronégatifs, soit à la composante associée aux individus supposés séropositifs :

\[
P(Z_i = 1) = \pi, \quad P(Z_i = 2) = 1 - \pi
\]

Conditionnellement à $Z_i$, on a :

\[
X_i \mid Z_i = k \sim \mathcal{N}(\mu_k, \sigma_k^2), \quad k = 1,2
\]

\textbf{Vraisemblance} \\


Soit \(\theta = (\pi_1,\pi_2,\mu_1,\mu_2,\sigma_1,\sigma_2)\) l'ensemble des paramètres du modèle. Pour un échantillon observé \(x = (x_1,\ldots,x_n)\).

Si les variables latentes \(z_i\) étaient observées, on pourrait écrire la vraisemblance complète à partir des densités \(f(x_i,z_i \mid \theta)\). Cependant, dans notre cas, les \(z_i\) ne sont pas observées : on ne sait pas directement si une observation \(x_i\) appartient à la première ou à la deuxième composante du mélange.

La vraisemblance observée est donc obtenue en sommant sur les deux classes possibles :

\[
f(x_i \mid \theta)
=
\sum_{k=1}^{2} f(x_i,z_i=k \mid \theta)
\]

Or, pour chaque composante \(k\), on a :

\[
f(x_i,z_i=k \mid \theta)
=
\pi_k \, \phi(x_i;\mu_k,\sigma_k),
\quad k=1,2.
\]

Ainsi :

\[
f(x_i \mid \theta)
=
\sum_{k=1}^{2}
\pi_k \, \phi(x_i;\mu_k,\sigma_k)
\]

La vraisemblance du modèle s'écrit alors comme le produit des densités de toutes les observations :

\[
L(\theta)
=
\prod_{i=1}^{n}
\sum_{k=1}^{2}
\pi_k \, \phi(x_i;\mu_k,\sigma_k)
\]

Dans le cas de deux composantes, avec \(\pi_1=\pi\) et \(\pi_2=1-\pi\), on obtient :

\[
L(\theta)
=
\prod_{i=1}^{n}
\left[
\pi \, \phi(x_i;\mu_1,\sigma_1)
+
(1-\pi) \, \phi(x_i;\mu_2,\sigma_2)
\right]
\]

Cette fonction est difficile à maximiser directement car les classes d'appartenance des observations sont inconnues et la vraisemblance contient une somme à l'intérieur du produit. L'algorithme EM, introduit par Dempster, Laird et Rubin (1977), est alors utilisé pour estimer les paramètres du modèle en présence de variables latentes.

\subsubsection{Algorithme EM}

L’estimation des paramètres du modèle de mélange gaussien est réalisée à l’aide de l’algorithme EM (\textit{Expectation-Maximization}). Cet algorithme est utilisé lorsque les données proviennent de plusieurs populations sous-jacentes mais que l’appartenance exacte de chaque observation à une population n’est pas connue.

Dans notre cas, le modèle suppose l’existence de deux populations :
\begin{itemize}
    \item une population séronégative ;
    \item une population séropositive.
\end{itemize}

L’algorithme EM alterne alors entre deux étapes principales. \\

\textbf{Étape E (Expectation)}

Lors de cette étape, on estime pour chaque observation la probabilité d’appartenir à chacune des deux composantes du mélange gaussien.

Autrement dit, l’algorithme évalue si une observation a davantage de chances d’appartenir à la population négative ou positive en fonction des paramètres estimés à l’itération précédente.

% Ces probabilités sont appelées \textit{responsabilités}. \\

\textbf{Étape M (Maximization)}

À partir des probabilités calculées lors de l’étape E, les paramètres du modèle sont mis à jour :
\begin{itemize}
    \item proportions des composantes ;
    \item moyennes ;
    \item variances.
\end{itemize}

Chaque observation contribue alors à l’estimation des paramètres de manière pondérée selon sa probabilité d’appartenance à chaque composante.

Les étapes E et M sont répétées successivement jusqu’à convergence de l’algorithme. \\

\textbf{Classification des individus} \\

Une fois le modèle ajusté, chaque observation est associée à une probabilité d’appartenance à chaque classe du modèle, c’est-à-dire à la classe 1 ou à la classe 2 dans notre cas. Ces deux classes sont ensuite interprétées en fonction des valeurs sérologiques estimées : la classe présentant les valeurs les plus faibles correspond aux individus supposés séronégatifs, tandis que la classe présentant les valeurs les plus élevées correspond aux individus supposés séropositifs. Ainsi, chaque individu possède une probabilité d’appartenir à la classe des séronégatifs et une probabilité d’appartenir à la classe des séropositifs.

On assigne ensuite chaque individu à la classe pour laquelle sa probabilité d'appartenance est la plus élevée. Par exemple, si un individu a une probabilité plus forte d'appartenir à la classe associée aux valeurs sérologiques élevées, il sera classé comme supposé séropositif. À l'inverse, si sa probabilité est plus forte pour la classe associée aux valeurs faibles, il sera classé comme supposé séronégatif.

Le tableau suivant donne quelques exemples fictifs de classification à partir des probabilités d'appartenance aux deux classes :

\begin{table}[h!]
\centering
\begin{tabular}{|c|c|c|c|}
\hline
Individu & Probabilité classe séronégative & Probabilité classe séropositive & Classification retenue \\
\hline
1 & 0{,}92 & 0{,}08 & Séronégatif \\
2 & 0{,}35 & 0{,}65 & Séropositif \\
3 & 0{,}51 & 0{,}49 & Séronégatif, mais incertain \\
\hline
\end{tabular}
\caption{Exemples fictifs de classification d'individus à partir des probabilités d'appartenance aux deux classes du modèle.}
\end{table}

\subsubsection{Calcul du cut-off}

Une fois les paramètres du modèle de mélange gaussien estimés, il est nécessaire de définir un seuil permettant de séparer les individus séronégatifs et séropositifs.

Dans ce projet, deux approches ont été utilisées pour déterminer ce seuil. \\

\textbf{Méthode 1 : intersection des distributions (\texttt{uniroot})}

La première approche consiste à déterminer le point d’intersection entre les deux composantes gaussiennes du mélange.

Le seuil correspond alors à la valeur $c$ telle que :

\[
\pi_1 \phi(c ; \mu_1, \sigma_1)
=
\pi_2 \phi(c ; \mu_2, \sigma_2)
\]

où :
\begin{itemize}
    \item $\pi_1$ et $\pi_2$ représentent les proportions des deux composantes ;
    \item $\phi$ désigne la densité d’une loi normale ;
    \item $(\mu_1, \sigma_1)$ et $(\mu_2, \sigma_2)$ sont les paramètres des deux distributions gaussiennes.
\end{itemize}

Le seuil est obtenu numériquement à l’aide de la fonction \texttt{uniroot} de \texttt{R}, qui permet de résoudre cette équation.

Cette méthode correspond au point où une observation a autant de chances d’appartenir à la population négative qu’à la population positive. \\

\begin{figure}[H]
\centering
\begin{tikzpicture}
\begin{axis}[
    width=11cm,
    height=6cm,
    axis lines=left,
    xlabel={Valeur standardisée},
    ylabel={Densité},
    xmin=-4, xmax=6,
    ymin=0, ymax=0.45,
    xtick=\empty,
    ytick=\empty,
    xlabel style={at={(axis description cs:0.5,-0.08)},anchor=north},
    ylabel style={at={(axis description cs:-0.06,0.5)},anchor=south},
    clip=false
]

% Paramètres visuels
\def\muone{0}
\def\mutwo{3.4}
\def\sigmaone{1.15}
\def\sigmatwo{1.05}
\def\cutoff{1.75}

% Distribution négative
\addplot[
    domain=-4:5,
    samples=300,
    blue,
    thick
]
{1/(\sigmaone*sqrt(2*pi))*exp(-((x-\muone)^2)/(2*\sigmaone^2))};

% Distribution positive
\addplot[
    domain=-2:6,
    samples=300,
    red,
    thick,
    dashed
]
{1/(\sigmatwo*sqrt(2*pi))*exp(-((x-\mutwo)^2)/(2*\sigmatwo^2))};


% Ligne cut-off
\addplot[
    black,
    dashed,
    thick
]
coordinates {(\cutoff,0) (\cutoff,0.42)};

% Labels
\node[black] at (axis cs:\cutoff,-0.035) {$c$};

\node[black] at (axis cs:\cutoff,0.45) {cut-off};
\node[black] at (axis cs:\cutoff,0.415) {uniroot};

\end{axis}
\end{tikzpicture}
\caption{Illustration de la méthode \texttt{uniroot}. La courbe bleue représente la population négative et la courbe rouge la population positive. Le seuil $c$ correspond au point d'intersection entre les deux composantes gaussiennes estimées par le modèle de mélange.}
\label{fig:cutoff_uniroot}
\end{figure}

\textbf{Méthode 2 : approche de type outlier sur la composante négative}

Une seconde approche consiste à considérer la première composante du mélange gaussien comme représentant la population séronégative.

Le seuil est alors défini à partir des paramètres de cette composante selon :

\[
\text{cut-off}
=
\mu_1 + 3\sigma_1
\]

où $\mu_1$ et $\sigma_1$ correspondent respectivement à la moyenne et à l’écart-type de la composante négative estimée par le GMM.

Cette approche est similaire à la méthode des valeurs extrêmes, mais utilise ici les paramètres estimés par le modèle de mélange gaussien.


\subsubsection{Estimation de la prévalence}

Une fois les seuils déterminés, la prévalence est estimée comme la proportion d’individus classés comme séropositifs.

Pour chacune des méthodes de détermination du cut-off, un individu est considéré comme positif lorsque sa valeur transformée dépasse le seuil estimé. \\

\newpage 

\textbf{Prévalence associée à la méthode uniroot (a) }

Pour la méthode basée sur l’intersection des distributions gaussiennes, la prévalence est définie par :

\[
p_{a}
=
\frac{1}{n}
\sum_{i=1}^{n}
\mathbf{1}_{(X_i > c_a)}
\]

où $c_a$ représente le seuil obtenu par intersection des deux composantes du mélange. \\

\textbf{Prévalence associée à la méthode outlier (b)} 

Pour l’approche basée sur la composante négative du GMM, la prévalence est estimée par :

\[
p_b
=
\frac{1}{n}
\sum_{i=1}^{n}
\mathbf{1}_{(X_i > c_b)}
\]

où :

\[
c_b
=
\mu_1 + 3\sigma_1
\]

avec $\mu_1$ et $\sigma_1$ les paramètres de la composante négative estimée par le modèle de mélange gaussien.


\textbf{Avantages et limites} \\

Les modèles de mélanges gaussiens présentent plusieurs avantages :

\begin{itemize}
    \item modélisation explicite de plusieurs populations ;
    \item estimation probabiliste du statut sérologique ;
    \item robustesse en présence de chevauchement entre distributions.
\end{itemize}

Cependant, ces modèles présentent également certaines limites :

\begin{itemize}
    \item sensibilité à l’initialisation ;
    \item hypothèse de normalité des composantes ;
    \item complexité plus élevée que la méthode de cut-off.
\end{itemize}

\section{Partie pratique}

\textbf{Environnement logiciel}

L’ensemble des analyses statistiques et des visualisations réalisées dans ce projet ont été effectués avec le logiciel \texttt{R} (version 4.5.2) à l’aide de l’environnement de développement \texttt{RStudio}. 

\newpage 

Plusieurs packages ont été utilisés afin de réaliser les différentes étapes de l’analyse :

\begin{itemize}
    \item \texttt{tidyverse} : manipulation et transformation des données ;
    \item \texttt{ggplot2} : réalisation des graphiques et visualisations ;
    \item \texttt{mclust} : ajustement des modèles de mélanges gaussiens (GMM) ;
    \item \texttt{gt} et \texttt{gtsummary} : création et mise en forme des tableaux ;
    \item \texttt{dplyr} et \texttt{tidyr} : restructuration et traitement des données ;
    \item \texttt{knitr} : génération des tableaux au format \LaTeX.
\end{itemize}

Les analyses ont été réalisées sur les valeurs transformées selon la relation :

\[
\log_{10}(MFI + 1)
\]

afin de réduire l’asymétrie des distributions et de faciliter l’ajustement des modèles statistiques.

\subsection{Jeu de données réel}

Le tableau suivant présente les principales caractéristiques descriptives du jeu de données utilisé dans cette étude. Les variables considérées incluent le pays d’origine des individus, le sexe, la classe d’âge ainsi que certaines informations biologiques comme le statut de gestation et de lactation. 

\begin{table}[H]
\caption{Tableau descriptif}
\centering
\begin{tabular}{|c|l|}
\hline
\multicolumn{1}{|l|}{\textbf{Characteristic}} & \textbf{N = 9\,593} \\ \hline
\multicolumn{1}{|l|}{\textbf{Pays}} & n (\%) \\ \hline
Cameroun & 4,885 (51) \\
RDC & 1,764 (18) \\
Republique de Guinee & 2,944 (31) \\ \hline
\multicolumn{1}{|l|}{\textbf{Sexe}} & \\ \hline
Femelle & 5,079 (54) \\
Male & 4,408 (46) \\
Inconnu & 106 \\ \hline
\multicolumn{1}{|l|}{\textbf{Age}} & \\ \hline
Adulte & 7,435 (82) \\
Immature & 437 (4.8) \\
Juvenile & 1,186 (13) \\
Inconnu  & 535 \\ \hline
\multicolumn{1}{|l|}{\textbf{Gestante}} & \\ \hline
Non & 7,771 (96) \\
Oui & 346 (4.3) \\
Inconnu & 1,476 \\ \hline
\multicolumn{1}{|l|}{\textbf{Lactation}} & \\ \hline
Non & 8,830 (96) \\
Oui & 413 (4.5) \\
Inconnu & 350 \\ \hline
\end{tabular}
\end{table}

Le jeu de données comprend un total de 9\,593 individus provenant principalement du Cameroun (51 \%), de la République de Guinée (31 \%) et de la RDC (18 \%).

La répartition selon le sexe est relativement équilibrée, avec une légère majorité de femelles (54 \%). Concernant l’âge, la majorité des individus sont classés comme adultes (82 \%), tandis que les individus immatures et juvéniles représentent une proportion plus faible.

Les variables de gestation et de lactation présentent majoritairement des réponses négatives, avec respectivement 96 \% des individus classés comme non gestants et non lactants.

Certaines observations apparaissent cependant comme \textit{Inconnu}. Ces valeurs correspondent à des informations manquantes ou non renseignées lors de la collecte des données sur le terrain. Cela peut s’expliquer par des difficultés d’identification biologique, des données incomplètes ou des contraintes liées aux conditions de prélèvement. 

Ces variables descriptives permettent de mieux caractériser la population étudiée avant l’analyse des distributions de MFI. 

On va se concentrer sur l'étude de 11 antigènes : 

\begin{itemize}
\item "GP.BEBV" , "GP.BOMV" , "GP.REBV" , "GP.SEBV", "GP.ZEBVkiss", "GP.ZEBVmay" 
\item"NP.SEBV" , "NP.ZEBV" 
\item "VP40.BEBV" , "VP40.SEBV" , "VP40.ZEBV" \\
\end{itemize}


La figure suivante présente les histogrammes et les distributions des valeurs transformées en log10(MFI + 1) pour les différents antigènes étudiés : 

\begin{figure}[H]
    \centering
    \includegraphics[width=0.8\textwidth]{figures/graphe_mfi.pdf}
    \caption{Histogramme et densité des MFI}
    \label{}
\end{figure}

\textbf{Observation} : \\
On observe globalement des distributions fortement asymétriques à droite, ce qui justifie l’utilisation d’une transformation logarithmique.

\subsubsection{Méthode outlier}

La figure suivante illustre la méthode outlier appliquée aux différentes distributions : 

\begin{figure}[H]
    \centering
    \includegraphics[width=0.8\textwidth]{figures/graphe_outlier_mfi.pdf}
    \caption{Méthode outlier - distribution et cut-off}
    \label{}
\end{figure}

\textbf{Observation} : \\
Les courbes rouges correspondent aux distributions normales ajustées à partir des populations négatives, tandis que les lignes verticales représentent les cut-offs estimés. \\
On observe que les seuils sont généralement situés dans la queue droite des distributions, ce qui permet de distinguer les individus potentiellement séropositifs.

\subsubsection{Modèle de mélange Gaussien }

La figure suivante présente les modèles de mélanges gaussiens ajustés aux distributions de MFI :

\begin{figure}[H]
    \centering
    \includegraphics[width=0.8\textwidth]{figures/graphe_gmm_mfi.pdf}
    \caption{GMM - distributions et cut-offs}
    \label{}
\end{figure}

\textbf{Observation} : \\
Les deux courbes représentent les composantes gaussiennes estimées correspondant aux populations négatives en rouge et positives en bleu. \\
Les lignes verticales représentent les cut-offs estimés par les deux méthodes , uniroot en noir et outlier en rouge. \\
Certaines distributions présentent une séparation relativement nette entre les deux composantes, notamment pour GP.SEBV ou GP.ZEBVkiss, tandis que d’autres présentent un chevauchement important. \\

\subsection{Résultats}

\begin{table} [H]

\caption{Paramètres estimés par la méthode outlier et le GMM}
\centering
\begin{tabular}{|l|rc|rrrr|}
\hline
 & \multicolumn{2}{c|}{Outlier} & & \multicolumn{2}{c}{GMM} & \\ \hline
MFI & mean & sd & mean{[}1{]} & \multicolumn{1}{r|}{sd{[}1{]}} & mean{[}2{]} & sd{[}2{]} \\ \hline
GP.BEBV     & 0.26          & 0.24         & 0.40        & \multicolumn{1}{r|}{0.34}      & 1.49        & 0.65      \\
GP.BOMV     & 0.21          & 0.18         & NA          & \multicolumn{1}{r|}{NA}        & NA          & NA        \\
GP.REBV     & 0.11          & 0.14         & NA          & \multicolumn{1}{r|}{NA}        & NA          & NA        \\
GP.SEBV     & 0.37          & 0.30         & 0.45        & \multicolumn{1}{r|}{0.37}      & 1.94        & 0.79      \\
GP.ZEBVkiss & 0.26          & 0.23         & 0.38        & \multicolumn{1}{r|}{0.32}      & 1.59        & 0.71      \\
GP.ZEBVmay  & 0.26          & 0.25         & 0.37        & \multicolumn{1}{r|}{0.33}      & 1.63        & 0.72      \\
NP.SEBV     & 0.17          & 0.17         & 0.38        & \multicolumn{1}{r|}{0.32}      & 1.19        & 0.62      \\
NP.ZEBV     & 0.19          & 0.18         & 0.40        & \multicolumn{1}{r|}{0.32}      & 1.05        & 0.55      \\
VP40.BEBV   & 0.21          & 0.19         & 0.42        & \multicolumn{1}{r|}{0.33}      & 1.08        & 0.56      \\
VP40.SEBV   & 0.34          & 0.24         & 0.50        & \multicolumn{1}{r|}{0.35}      & 1.14        & 0.54      \\
VP40.ZEBV   & 0.24          & 0.20         & 0.48        & \multicolumn{1}{r|}{0.34}      & 1.18        & 0.56      \\ \hline
\end{tabular}
\end{table}

Le tableau 3 présente les moyennes et écarts-types estimés par les méthodes outlier et GMM.

On observe que les moyennes associées à la seconde composante du GMM sont systématiquement plus élevées que celles de la première composante, ce qui est cohérent avec l’interprétation des populations positives et négatives.

Pour certains antigènes, le modèle de mélange Gaussien n’a pas permis de distinguer correctement deux composantes gaussiennes, ce qui a conduit à des valeurs manquantes (NA). Cela suggère que certaines distributions ne présentent pas deux populations suffisamment séparées pour être modélisées correctement par un GMM à deux composantes.

\begin{table}[H]

\caption{Cut-offs et prévalences (en \%) estimées par antigène}
\centering
\begin{tabular}{|l|rr|rrrr|}
\hline
& \multicolumn{2}{c|}{Outlier} & & \multicolumn{2}{c}{GMM} & \\ \hline
 & & & \multicolumn{2}{c|}{uniroot} & \multicolumn{2}{c|}{outlier} \\ \hline
MFI & cut-off & prévalence & cut-off & \multicolumn{1}{r|}{prévalence} & cut-off  & prévalence \\ \hline
GP.BEBV     & 0.98        & 25.7           & 0.98    & \multicolumn{1}{r|}{32.5}       & 1.43     & 13.1       \\ \hline
GP.BOMV     & 0.75        & 8.6            & NA      & \multicolumn{1}{r|}{NA}         & 1.44     & 0.7        \\ \hline
GP.REBV     & 0.54        & 31.1           & NA      & \multicolumn{1}{r|}{NA}         & 1.87    & 0.5        \\ \hline
GP.SEBV     & 1.27        & 30.4           & 1.09    & \multicolumn{1}{r|}{42.0}       & 1.55     & 23.0       \\ \hline
GP.ZEBVkiss & 0.96        & 29.7           & 0.94    & \multicolumn{1}{r|}{36.4}       & 1.34     & 18.5       \\ \hline
GP.ZEBVmay  & 1.00        & 29.6           & 0.94    & \multicolumn{1}{r|}{37.8}       & 1.36     & 19.4       \\ \hline
NP.SEBV     & 0.69        & 20.7           & 0.99    & \multicolumn{1}{r|}{17.3}       & 1.33     & 5.1        \\ \hline
NP.ZEBV     & 0.73        & 17.2           & 0.98    & \multicolumn{1}{r|}{15.5}       & 1.35     & 2.8        \\ \hline
VP40.BEBV   & 0.76        & 19.2           & 1.02    & \multicolumn{1}{r|}{15.9}       & 1.40     & 2.9        \\ \hline
VP40.SEBV   & 1.05        & 15.9           & 0.97    & \multicolumn{1}{r|}{30.0}       & 1.54     & 3.4        \\ \hline
VP40.ZEBV   & 0.84        & 15.5           & 1.15    & \multicolumn{1}{r|}{11.4}       & 1.50     & 2.3        \\ \hline
\end{tabular}

\end{table}

Le tableau 4 présente les cut-offs et les prévalences estimées selon les différentes méthodes.

Globalement, les prévalences obtenues avec la méthode GMM uniroot sont souvent plus élevées que celles obtenues avec la méthode outlier. \\
En revanche, la méthode GMM outlier produit généralement des prévalences plus faibles en raison de seuils plus élevés.\\
Pour certains antigènes, le modèle GMM n’a pas permis d’estimer correctement les paramètres, conduisant à des valeurs manquantes (NA). Cela suggère l’absence de structure bimodale suffisamment marquée dans les distributions correspondantes.

\section{Conclusion (ouverture)}

\section{Bibliographie}

\bibliographystyle{unsrt}
\bibliography{references}

\section{Annexe}

L’ensemble des scripts \texttt{R}, des figures et des tableaux utilisés dans ce projet est disponible sur le dépôt GitHub suivant :

\begin{center}
\href{https://github.com/Bastien2003/Projet-M1SSD}{Dépôt GitHub du projet}
\end{center}

\end{document}